\chapter{Regular Surfaces}
\label{ch:vii14-regular-surfaces}

A regular surface is a two-dimensional smooth manifold realized locally inside \(\mathbb R^3\).  The central local object is a surface patch
\[
X(u,v):U\subset\mathbb R^2\longrightarrow\mathbb R^3
\]
whose two coordinate derivatives are linearly independent.  Equivalently,
\[
\boxed{X_u\times X_v\neq0.}
\]
This rank condition produces tangent planes, local normals, coordinate curves, and the framework needed for the fundamental forms.

\section{Regular surface patches}
\begin{definition}[Regular parametrized patch]\label{def:vii14-patch}
A smooth map \(X:U\to\mathbb R^3\), with \(U\subset\mathbb R^2\) open, is a regular patch if \(dX_{(u,v)}\) has rank \(2\) at every point and \(X\) is a homeomorphism onto its image.
\end{definition}

\begin{proposition}[Cross-product criterion]\label{prop:vii14-cross-criterion}
A smooth parametrization has rank \(2\) exactly when
\[
X_u\times X_v\neq0.
\]
\end{proposition}
\begin{proof}
The derivative has columns \(X_u,X_v\).  It has rank \(2\) exactly when these vectors are linearly independent, equivalent in \(\mathbb R^3\) to a nonzero cross product.
\end{proof}

\begin{definition}[Regular surface]\label{def:vii14-regular-surface}
A subset \(S\subset\mathbb R^3\) is a regular surface if every point \(p\in S\) lies in the image of a regular patch whose image is open in \(S\) with the subspace topology.
\end{definition}

\begin{warning}[Rank alone is not the whole patch definition]\label{warn:vii14-rank-topology}
The rank condition is differential.  A surface patch must also give the correct local topology; global self-overlap can prevent a parametrization from being a chart even if its derivative has rank \(2\).
\end{warning}

\section{Tangent planes}
\begin{definition}[Tangent plane]\label{def:vii14-tangent-plane}
For \(p=X(u_0,v_0)\), define
\[
T_pS=\operatorname{span}\{X_u(u_0,v_0),X_v(u_0,v_0)\}.
\]
\end{definition}

\begin{theorem}[Curve characterization]\label{thm:vii14-curve-tangent}
A vector \(w\in\mathbb R^3\) lies in \(T_pS\) if and only if there is a smooth curve \(\alpha\) in \(S\) with \(\alpha(0)=p\) and \(\alpha'(0)=w\).
\end{theorem}
\begin{proof}
If \(\alpha(t)=X(u(t),v(t))\), then \(\alpha'=u'X_u+v'X_v\).  Conversely any combination \(aX_u+bX_v\) is realized by the coordinate-domain line \((u_0+at,v_0+bt)\) for small \(t\).
\end{proof}

\begin{definition}[Local unit normal]\label{def:vii14-unit-normal}
On a regular patch,
\[
N=\frac{X_u\times X_v}{\|X_u\times X_v\|}
\]
is a smooth unit normal field.
\end{definition}

\begin{proposition}[Tangent-plane equation]\label{prop:vii14-plane-equation}
At \(p\) with unit normal \(N(p)\), the affine tangent plane is
\[
\{q\in\mathbb R^3:\langle q-p,N(p)\rangle=0\}.
\]
\end{proposition}

\section{Change of parameters}
\begin{theorem}[Reparametrization]\label{thm:vii14-reparam}
Let \(\phi:\widetilde U\to U\) be a diffeomorphism and \(\widetilde X=X\circ\phi\).  Then \(\widetilde X\) is regular whenever \(X\) is regular, and the tangent plane is unchanged.
\end{theorem}
\begin{proof}
The chain rule gives \(d\widetilde X=dX\circ d\phi\).  Since \(d\phi\) is invertible, both derivatives have the same rank and image.
\end{proof}

\begin{proposition}[Normal orientation under coordinates]\label{prop:vii14-normal-change}
Under a change of parameters, the vector \(X_u\times X_v\) is multiplied by the Jacobian determinant of the parameter change.  Thus orientation-preserving changes preserve the chosen normal direction, while orientation-reversing changes reverse it.
\end{proposition}

\section{Graphs and implicit surfaces}
\begin{proposition}[Graphs are regular]\label{prop:vii14-graph}
For smooth \(f:U\subset\mathbb R^2\to\mathbb R\),
\[
X(u,v)=(u,v,f(u,v))
\]
is a regular patch.  A normal vector is
\[
(-f_u,-f_v,1).
\]
\end{proposition}

\begin{theorem}[Implicit surface criterion]\label{thm:vii14-implicit}
Let \(F:\mathbb R^3\to\mathbb R\) be smooth.  If \(\nabla F(p)\neq0\) on \(F^{-1}(c)\), then the level set is locally a regular surface and
\[
T_pS=\ker dF_p=\nabla F(p)^\perp.
\]
\end{theorem}
\begin{proof}
This is the regular-level theorem of VII/06 in dimension \(3\to1\).  The implicit-function theorem also writes the level locally as a graph.
\end{proof}

\begin{example}[Sphere]\label{ex:vii14-sphere}
The sphere \(S_R^2\) is the regular level of \(F(x,y,z)=x^2+y^2+z^2\) at \(R^2\).  Hence
\[
T_pS_R^2=p^\perp.
\]
Spherical coordinates provide useful patches away from their polar coordinate singularities.
\end{example}

\begin{example}[Cylinder]\label{ex:vii14-cylinder}
The circular cylinder has patch
\[
X(u,v)=(R\cos u,R\sin u,v),
\]
with \(X_u\times X_v\neq0\) everywhere.
\end{example}

\section{Coordinate curves and surface velocities}
\begin{definition}[Coordinate curves]\label{def:vii14-coordinate-curves}
For a patch \(X(u,v)\), curves with \(v=v_0\) and with \(u=u_0\) are the coordinate curves.  Their tangent vectors are \(X_u\) and \(X_v\), respectively.
\end{definition}

\begin{proposition}[Velocity in surface coordinates]\label{prop:vii14-surface-velocity}
If \(\alpha(t)=X(u(t),v(t))\), then
\[
\alpha'=u'X_u+v'X_v.
\]
Thus every surface velocity decomposes uniquely in the coordinate tangent basis.
\end{proposition}

\section{Standard families of surfaces}
\begin{example}[Surface of revolution]\label{ex:vii14-revolution}
If a profile \((r(v),z(v))\) with \(r(v)>0\) is revolved around the \(z\)-axis,
\[
X(u,v)=(r(v)\cos u,r(v)\sin u,z(v)).
\]
Regularity requires the profile to be regular and avoids points where the rotational coordinate degenerates.
\end{example}

\begin{example}[Torus]\label{ex:vii14-torus}
For \(R>r>0\),
\[
X(u,v)=((R+r\cos v)\cos u,(R+r\cos v)\sin u,r\sin v)
\]
is regular.  Both parameters are periodic, so a global one-to-one rectangular patch is impossible, but finitely many patches cover the torus.
\end{example}

\begin{example}[Helicoid]\label{ex:vii14-helicoid}
The map
\[
X(u,v)=(u\cos v,u\sin v,av),\qquad a\neq0,
\]
defines a regular helicoid.  Its coordinate curves exhibit simultaneously radial and helical directions.
\end{example}

\section{Overlap maps and intrinsic smoothness}
\begin{theorem}[Smooth overlap maps]\label{thm:vii14-overlaps}
If \(X:U\to S\) and \(Y:V\to S\) are overlapping regular surface patches, then
\[
Y^{-1}\circ X
\]
is smooth wherever defined.
\end{theorem}
\begin{proof}[Proof idea]
Near an overlap point, choose a coordinate projection of \(\mathbb R^3\) whose restriction has nonsingular derivative on the tangent plane.  The inverse-function theorem makes that projection a local chart for both patches, and the transition is a composition of smooth local inverses.
\end{proof}

\begin{corollary}[Surface as a smooth manifold]\label{cor:vii14-surface-manifold}
Every regular surface in \(\mathbb R^3\) carries a canonical smooth \(2\)-manifold structure compatible with its regular patches.
\end{corollary}

\section{Exercises}
\begin{exercise}\label{exr:vii14-01}
Check that \(X(u,v)=(u,v,u^2+v^2)\) is regular.
\end{exercise}
\begin{hint}
% PEDAGOGY-ENRICHED-VII
Compute \(X_u\times X_v\). Make the controlling geometric criterion explicit before the calculation. Work in a regular surface chart or level-set description and verify rank before using tangent-plane or local-coordinate formulas.
\end{hint}
\begin{solution}
\(X_u=(1,0,2u)\), \(X_v=(0,1,2v)\), so their cross product \((-2u,-2v,1)\) never vanishes.
\end{solution}

\begin{exercise}\label{exr:vii14-02}
Find the tangent plane to \(z=x^2+y^2\) at \((1,0,1)\).
\end{exercise}
\begin{hint}
% PEDAGOGY-ENRICHED-VII
Use the graph normal. Work in a convenient local model first, then return to the invariant statement. Work in a regular surface chart or level-set description and verify rank before using tangent-plane or local-coordinate formulas.
\end{hint}
\begin{solution}
A normal is \((-2,0,1)\), so the plane is \(-2(x-1)+(z-1)=0\).
\end{solution}

\begin{exercise}\label{exr:vii14-03}
Prove the cross-product rank criterion.
\end{exercise}
\begin{hint}
% PEDAGOGY-ENRICHED-VII
Rank two means two independent columns. After the main computation, check the hypotheses and geometric interpretation. Work in a regular surface chart or level-set description and verify rank before using tangent-plane or local-coordinate formulas.
\end{hint}
\begin{solution}
The columns are \(X_u,X_v\); in \(\mathbb R^3\) they are independent exactly when their cross product is nonzero.
\end{solution}

\begin{exercise}\label{exr:vii14-04}
For the sphere show \(T_pS^2=p^\perp\).
\end{exercise}
\begin{hint}
% PEDAGOGY-ENRICHED-VII
Use the level-set function \(F(x)=\|x\|^2\). Make the controlling geometric criterion explicit before the calculation. Work in a regular surface chart or level-set description and verify rank before using tangent-plane or local-coordinate formulas.
\end{hint}
\begin{solution}
\(dF_p(v)=2p\cdot v\), so the tangent space is its kernel \(p^\perp\).
\end{solution}

\begin{exercise}\label{exr:vii14-05}
Check the cylinder patch is regular.
\end{exercise}
\begin{hint}
% PEDAGOGY-ENRICHED-VII
Compute its coordinate derivatives. Work in a convenient local model first, then return to the invariant statement. Work in a regular surface chart or level-set description and verify rank before using tangent-plane or local-coordinate formulas.
\end{hint}
\begin{solution}
\(X_u=(-R\sin u,R\cos u,0)\), \(X_v=(0,0,1)\); their cross product has norm \(R\).
\end{solution}

\begin{exercise}\label{exr:vii14-06}
Why do spherical coordinates fail to be regular at a pole?
\end{exercise}
\begin{hint}
% PEDAGOGY-ENRICHED-VII
Inspect \(X_\theta\times X_\phi\). After the main computation, check the hypotheses and geometric interpretation. Work in a regular surface chart or level-set description and verify rank before using tangent-plane or local-coordinate formulas.
\end{hint}
\begin{solution}
One coordinate vector collapses at a pole, so the cross product vanishes even though the sphere itself is regular there.
\end{solution}

\begin{exercise}\label{exr:vii14-07}
Prove graph patches are regular.
\end{exercise}
\begin{hint}
% PEDAGOGY-ENRICHED-VII
Look at the first two components of \(X_u,X_v\). Make the controlling geometric criterion explicit before the calculation. Work in a regular surface chart or level-set description and verify rank before using tangent-plane or local-coordinate formulas.
\end{hint}
\begin{solution}
They contain the independent vectors \((1,0)\) and \((0,1)\), so the full vectors are independent.
\end{solution}

\begin{exercise}\label{exr:vii14-08}
Find a unit normal to a graph.
\end{exercise}
\begin{hint}
% PEDAGOGY-ENRICHED-VII
Normalize \((-f_u,-f_v,1)\). Work in a convenient local model first, then return to the invariant statement. Work in a regular surface chart or level-set description and verify rank before using tangent-plane or local-coordinate formulas.
\end{hint}
\begin{solution}
\(N=(-f_u,-f_v,1)/\sqrt{1+f_u^2+f_v^2}\).
\end{solution}

\begin{exercise}\label{exr:vii14-09}
Prove the velocity formula \(\alpha'=u'X_u+v'X_v\).
\end{exercise}
\begin{hint}
% PEDAGOGY-ENRICHED-VII
Apply the chain rule. After the main computation, check the hypotheses and geometric interpretation. Work in a regular surface chart or level-set description and verify rank before using tangent-plane or local-coordinate formulas.
\end{hint}
\begin{solution}
Differentiate \(X(u(t),v(t))\) componentwise to obtain the formula.
\end{solution}

\begin{exercise}\label{exr:vii14-10}
Why is the tangent basis decomposition unique?
\end{exercise}
\begin{hint}
% PEDAGOGY-ENRICHED-VII
Use regularity. Make the controlling geometric criterion explicit before the calculation. Work in a regular surface chart or level-set description and verify rank before using tangent-plane or local-coordinate formulas.
\end{hint}
\begin{solution}
\(X_u,X_v\) are linearly independent, so coordinates in their span are unique.
\end{solution}

\begin{exercise}\label{exr:vii14-11}
Show a regular change of parameters preserves the tangent plane.
\end{exercise}
\begin{hint}
% PEDAGOGY-ENRICHED-VII
Use \(d(X\circ\phi)=dX\,d\phi\). Work in a convenient local model first, then return to the invariant statement. Work in a regular surface chart or level-set description and verify rank before using tangent-plane or local-coordinate formulas.
\end{hint}
\begin{solution}
Because \(d\phi\) is invertible, the image of the composite derivative equals the image of \(dX\).
\end{solution}

\begin{exercise}\label{exr:vii14-12}
What happens to the local normal under a parameter swap \((u,v)\mapsto(v,u)\)?
\end{exercise}
\begin{hint}
% PEDAGOGY-ENRICHED-VII
The Jacobian determinant is \(-1\). After the main computation, check the hypotheses and geometric interpretation. Work in a regular surface chart or level-set description and verify rank before using tangent-plane or local-coordinate formulas.
\end{hint}
\begin{solution}
The cross product reverses sign, so the chosen unit normal reverses.
\end{solution}

\begin{exercise}\label{exr:vii14-13}
Show \(x^2+y^2-z=0\) is a regular surface.
\end{exercise}
\begin{hint}
% PEDAGOGY-ENRICHED-VII
Check the gradient. Make the controlling geometric criterion explicit before the calculation. Work in a regular surface chart or level-set description and verify rank before using tangent-plane or local-coordinate formulas.
\end{hint}
\begin{solution}
\(\nabla F=(2x,2y,-1)\) never vanishes, so every level point is regular.
\end{solution}

\begin{exercise}\label{exr:vii14-14}
Where does \(x^2+y^2-z^2=0\) fail the implicit regularity test?
\end{exercise}
\begin{hint}
% PEDAGOGY-ENRICHED-VII
Compute \(\nabla F\). Work in a convenient local model first, then return to the invariant statement. Work in a regular surface chart or level-set description and verify rank before using tangent-plane or local-coordinate formulas.
\end{hint}
\begin{solution}
The gradient \((2x,2y,-2z)\) vanishes at the origin, the cone vertex.
\end{solution}

\begin{exercise}\label{exr:vii14-15}
Write the affine tangent-plane equation for an implicit surface.
\end{exercise}
\begin{hint}
% PEDAGOGY-ENRICHED-VII
Use the gradient as normal. After the main computation, check the hypotheses and geometric interpretation. Work in a regular surface chart or level-set description and verify rank before using tangent-plane or local-coordinate formulas.
\end{hint}
\begin{solution}
At \(p\), \(\nabla F(p)\cdot(q-p)=0\).
\end{solution}

\begin{exercise}\label{exr:vii14-16}
Check regularity of the helicoid patch.
\end{exercise}
\begin{hint}
% PEDAGOGY-ENRICHED-VII
Compute the cross product. Make the controlling geometric criterion explicit before the calculation. Work in a regular surface chart or level-set description and verify rank before using tangent-plane or local-coordinate formulas.
\end{hint}
\begin{solution}
\(X_u=(\cos v,\sin v,0)\), \(X_v=(-u\sin v,u\cos v,a)\); the cross product has squared norm \(a^2+u^2>0\).
\end{solution}

\begin{exercise}\label{exr:vii14-17}
Why can a torus not be covered by one periodic rectangle as a one-to-one patch?
\end{exercise}
\begin{hint}
% PEDAGOGY-ENRICHED-VII
Opposite parameter edges represent the same points. Work in a convenient local model first, then return to the invariant statement. Work in a regular surface chart or level-set description and verify rank before using tangent-plane or local-coordinate formulas.
\end{hint}
\begin{solution}
Periodicity forces identifications, violating global injectivity on a full closed-period rectangle; overlapping open patches are used instead.
\end{solution}

\begin{exercise}\label{exr:vii14-18}
Describe the coordinate curves on a cylinder.
\end{exercise}
\begin{hint}
% PEDAGOGY-ENRICHED-VII
Hold one parameter fixed. After the main computation, check the hypotheses and geometric interpretation. Work in a regular surface chart or level-set description and verify rank before using tangent-plane or local-coordinate formulas.
\end{hint}
\begin{solution}
Fixing height gives circles; fixing angle gives vertical lines.
\end{solution}

\begin{exercise}\label{exr:vii14-19}
For a surface of revolution compute \(X_u\).
\end{exercise}
\begin{hint}
% PEDAGOGY-ENRICHED-VII
Differentiate in the angular coordinate. Make the controlling geometric criterion explicit before the calculation. Work in a regular surface chart or level-set description and verify rank before using tangent-plane or local-coordinate formulas.
\end{hint}
\begin{solution}
\(X_u=(-r\sin u,r\cos u,0)\).
\end{solution}

\begin{exercise}\label{exr:vii14-20}
State a regularity condition for a surface of revolution away from the axis.
\end{exercise}
\begin{hint}
% PEDAGOGY-ENRICHED-VII
Inspect \(X_u\times X_v\). Work in a convenient local model first, then return to the invariant statement. Work in a regular surface chart or level-set description and verify rank before using tangent-plane or local-coordinate formulas.
\end{hint}
\begin{solution}
It is nonzero when \(r>0\) and the profile derivative \((r',z')\) is nonzero.
\end{solution}

\begin{exercise}\label{exr:vii14-21}
Explain why a regular surface is locally a graph after a coordinate permutation.
\end{exercise}
\begin{hint}
% PEDAGOGY-ENRICHED-VII
Use a nonzero component of the normal. After the main computation, check the hypotheses and geometric interpretation. Work in a regular surface chart or level-set description and verify rank before using tangent-plane or local-coordinate formulas.
\end{hint}
\begin{solution}
Rank two means some \(2\times2\) minor of \(dX\) is nonsingular; projection to the corresponding coordinate plane is locally invertible, yielding a graph.
\end{solution}

\begin{exercise}\label{exr:vii14-22}
Show the tangent plane is independent of the chosen surface curve realizing a tangent vector.
\end{exercise}
\begin{hint}
% PEDAGOGY-ENRICHED-VII
Use the patch derivative image. Make the controlling geometric criterion explicit before the calculation. Work in a regular surface chart or level-set description and verify rank before using tangent-plane or local-coordinate formulas.
\end{hint}
\begin{solution}
Every curve velocity lies in \(\operatorname{im}dX\), and every vector in that image is realized by a parameter-domain curve, so the plane is intrinsic.
\end{solution}

\begin{exercise}\label{exr:vii14-23}
What is the dimension of a regular surface tangent space?
\end{exercise}
\begin{hint}
% PEDAGOGY-ENRICHED-VII
Use rank two. Work in a convenient local model first, then return to the invariant statement. Work in a regular surface chart or level-set description and verify rank before using tangent-plane or local-coordinate formulas.
\end{hint}
\begin{solution}
It is two-dimensional because it is the image of a rank-two derivative.
\end{solution}

\begin{exercise}\label{exr:vii14-24}
Explain why the cone vertex is not a regular-surface point.
\end{exercise}
\begin{hint}
% PEDAGOGY-ENRICHED-VII
Compare local topology or tangent behavior. After the main computation, check the hypotheses and geometric interpretation. Work in a regular surface chart or level-set description and verify rank before using tangent-plane or local-coordinate formulas.
\end{hint}
\begin{solution}
No unique tangent plane exists there; equivalently the standard implicit gradient vanishes and the cone cannot be represented as a smooth graph near the vertex.
\end{solution}

\section{Solved problem dossiers}
\begin{problem}[Legacy regular-patch problem]\label{prob:vii14-01}
Prove the rank-two and cross-product criteria are equivalent.
\end{problem}
\begin{solution}
The derivative has the two coordinate tangent vectors as columns; rank two is their linear independence, equivalent to nonzero cross product.
\end{solution}

\begin{problem}[Legacy tangent-plane problem]\label{prob:vii14-02}
Prove the tangent plane equals the set of velocities of curves on the surface.
\end{problem}
\begin{solution}
Apply the chain rule in one direction and realize any coordinate linear combination by a straight parameter-domain curve in the other.
\end{solution}

\begin{problem}[Legacy graph problem]\label{prob:vii14-03}
Show a smooth graph is a regular surface and identify its normal.
\end{problem}
\begin{solution}
Its coordinate tangent vectors are independent, and their cross product is \((-f_u,-f_v,1)\).
\end{solution}

\begin{problem}[Legacy implicit-surface problem]\label{prob:vii14-04}
Recover a regular surface from a regular level set.
\end{problem}
\begin{solution}
If \(\nabla F\neq0\), the regular-level theorem gives a two-dimensional embedded submanifold with tangent space \(\ker dF=\nabla F^\perp\).
\end{solution}

\begin{problem}[Sphere atlas dossier]\label{prob:vii14-05}
Explain why the sphere is regular despite singularities of spherical coordinates.
\end{problem}
\begin{solution}
Regularity is a property of the surface, not one chosen parametrization. Near each point use an implicit-function graph or a coordinate projection patch with nonsingular derivative.
\end{solution}

\begin{problem}[Cylinder dossier]\label{prob:vii14-06}
Compute tangent basis and normal for the circular cylinder.
\end{problem}
\begin{solution}
The standard patch has \(X_u=(-R\sin u,R\cos u,0)\), \(X_v=(0,0,1)\), and a radial normal \((\cos u,\sin u,0)\).
\end{solution}

\begin{problem}[Change-of-parameters dossier]\label{prob:vii14-07}
Prove tangent planes and regularity are invariant under a diffeomorphic parameter change.
\end{problem}
\begin{solution}
The chain rule composes \(dX\) with an invertible \(2\times2\) derivative, preserving rank and image.
\end{solution}

\begin{problem}[Surface-of-revolution dossier]\label{prob:vii14-08}
Derive the regularity condition for a standard surface of revolution.
\end{problem}
\begin{solution}
Compute \(X_u\times X_v\); its norm is \(r\sqrt{r'^2+z'^2}\), so it is regular wherever \(r>0\) and the profile is regular.
\end{solution}

\begin{problem}[Torus dossier]\label{prob:vii14-09}
Verify the standard torus parametrization is regular for \(R>r>0\).
\end{problem}
\begin{solution}
The cross-product norm is \(r(R+r\cos v)\), strictly positive because \(R-r>0\).
\end{solution}

\begin{problem}[Helicoid dossier]\label{prob:vii14-10}
Verify regularity and identify tangent directions of the helicoid.
\end{problem}
\begin{solution}
The cross product has norm \(\sqrt{a^2+u^2}\). The \(u\)-curves are horizontal radial lines and the \(v\)-curves are helices.
\end{solution}

\begin{problem}[Overlap-map dossier]\label{prob:vii14-11}
Why are transition maps between regular surface patches smooth?
\end{problem}
\begin{solution}
Choose an ambient coordinate projection that is nonsingular on the tangent plane. It supplies a common local Euclidean chart, reducing the transition to compositions of smooth maps and inverse-function-theorem inverses.
\end{solution}

\begin{problem}[Singular parametrization versus singular surface]\label{prob:vii14-12}
Distinguish a bad parametrization from a genuinely singular surface point.
\end{problem}
\begin{solution}
A bad parametrization can lose rank even when another patch is regular, as at spherical-coordinate poles. A genuinely singular point admits no regular patch at all, as at the cone vertex.
\end{solution}

\section{Challenge problems}
\begin{problem}[Local graph theorem for surfaces]\label{prob:vii14-ch01}
Prove every regular surface is locally a graph over some coordinate plane.
\end{problem}
\begin{solution}
At a regular point some \(2\times2\) minor of the patch derivative is nonsingular. Project to those two ambient coordinates. The inverse-function theorem makes the projection composed with the patch a local diffeomorphism; solving for the remaining coordinate gives the graph representation.
\end{solution}

\begin{problem}[Equivalent tangent definitions]\label{prob:vii14-ch02}
Show that the derivative-image, curve-velocity, and implicit-kernel definitions of the tangent plane agree whenever all are available.
\end{problem}
\begin{solution}
The derivative-image and curve-velocity descriptions agree by the chain rule and curve realization. For a regular level \(F=c\), differentiating \(F\circ\alpha=c\) gives inclusion in \(\ker dF\), and both spaces have dimension two, hence equality.
\end{solution}

\begin{problem}[Atlas from regular patches]\label{prob:vii14-ch03}
Build explicitly the smooth-manifold atlas carried by a regular surface.
\end{problem}
\begin{solution}
Use inverses of regular patches as charts on their images. The overlap theorem gives smooth transitions. The subspace topology is second countable and Hausdorff as a subspace of \(\mathbb R^3\), so the resulting atlas defines a smooth two-manifold.
\end{solution}

\begin{problem}[Orientability preview]\label{prob:vii14-ch04}
Explain how local normals on overlapping patches encode the obstruction to a global orientation.
\end{problem}
\begin{solution}
Each patch supplies a unit normal. On a connected overlap two choices differ by sign. An orientation is a coherent choice with positive sign on all overlaps; failure to make these signs globally consistent is precisely the obstruction developed later in surface orientation theory.
\end{solution}

\begin{problem}[Regularity under ambient diffeomorphism]\label{prob:vii14-ch05}
If \(\Phi:\mathbb R^3\to\mathbb R^3\) is a diffeomorphism and \(S\) is regular, prove \(\Phi(S)\) is regular.
\end{problem}
\begin{solution}
Compose every regular patch \(X\) with \(\Phi\). Its derivative is \(d\Phi\circ dX\), rank two because \(d\Phi\) is invertible, and the composition is a homeomorphism onto the corresponding image.
\end{solution}

\section*{Bridge to VII/15}
A regular patch provides a tangent basis \(X_u,X_v\) and a local normal \(N\).  VII/15 uses these to measure lengths, angles, areas, and normal bending through the first and second fundamental forms.
